\chapter{Options Market Making in Practice}\label{ch:dv:options-market-making-in-practice}

On the eve of an ex-dividend date, two market makers trade 33\,000 contracts of a deep in-the-money call with each other, in both directions.
Neither takes a position. Both exercise every call they bought. The next morning each finds that some of the calls it sold were never assigned:
their holders, among the public, failed to exercise. Each unassigned short call is worth the dividend less a little time value. In this chapter's
model, with 30\% of 10\,000 public contracts left unexercised, the pair clears \$108\,553 after fees. The \omterm{def:dv:options-market-making-in-practice:play}{dividend play} is a small corner of options
market making, but it shows the trade's character: many small, well-defined edges, each engineered, each with its own costs and risks. This
chapter builds the rest of them: a live surface and \omterm{def:dv:options-market-making-in-practice:theo}{theoretical values}, quote widths, delta-hedging policies, dividends and pin risk, and the
limits that keep a book's risk in buckets.

\section{Fitting a live surface}

A market maker quotes thousands of series. It cannot price each from scratch as the underlying moves, so it keeps a surface: a small set of
parameters per expiry (chapter 8's SVI, or a skew model), refitted to the market's quotes and trades at intervals, and moved with the spot in
between by a marking rule (chapter 7).

\begin{definition}[Theoretical value]\label{def:dv:options-market-making-in-practice:theo}
The \emph{theoretical value}\index{theoretical value} of an option, for a market maker, is its price on the market maker's current surface,
forward and discount curve: the reference around which it sets its bid and ask, updated with every move of the underlying and every refit.
\end{definition}

Between refits the marking rule does the work. If the market's smile moves with the spot (\omterm{def:dv:implied-volatility-and-its-surface:sticky}{sticky delta}) but the theo holds each strike's volatility
(\omterm{def:dv:implied-volatility-and-its-surface:sticky}{sticky strike}), the theo drifts off by the skew times the log move of the spot since the fit, divided by $\sqrt\tau$. With a skew of $-0.10$ and a
spot at 20\% volatility, the mean error in volatility points is:

\begin{center}\small
\begin{tabular}{lrrrr}
\toprule
refit interval & 1 second & 10 seconds & 1 minute & 5 minutes\\
\midrule
one-month options & 0.0023 & 0.0072 & 0.018 & 0.039\\
three-month options & 0.0013 & 0.0042 & 0.010 & 0.023\\
\bottomrule
\end{tabular}
\end{center}

A minute's staleness costs a fiftieth of a volatility point on a one-month option, well inside any width. What makes surfaces go wrong is not the
interval but the jumps: an earnings release, a large trade or an index move shifts the whole smile at once. The quoting engine therefore refits on
events as well as on a clock, and pulls its quotes while it refits.

\section{Theoretical value and quote widths}

Around the theo the market maker quotes a width. Too narrow, and informed traders pick it off whenever its theo is wrong; too wide, and the
uninformed flow that pays for everything goes elsewhere. Book 2, chapter 30 treated the width in a game; the build turns it into a formula per
series.

\begin{example}[The width model]\label{ex:dv:options-market-making-in-practice:width}
Measure everything in volatility points. The theo's error is normal with a standard deviation of 0.5 point. Uninformed orders arrive at 20 an hour
when the half-width is zero, falling as $e^{-w/0.3}$ as it widens; each pays the half-width $w$. Informed orders arrive at $\lambda$ an hour and
trade only when the error exceeds $w$, taking the excess. The expected profit per hour, per unit of \omterm{def:dv:greeks-and-the-hedging-pnl:greeks}{vega}, is
\[
20\,e^{-w/0.3}\,w-2\lambda\bigl(0.5\,\varphi(w/0.5)-w\,(1-\Phi(w/0.5))\bigr).
\]
Without informed flow the best half-width is 0.30 point and earns 2.21 an hour. With 2, 5, 10 and 20 informed orders an hour it widens to 0.35,
0.42, 0.55 and 0.78, and the profit falls to 1.89, 1.51, 1.07 and 0.65
(\cref{fig:dv:options-market-making-in-practice:width}).
\end{example}

\begin{omfigure}
\begin{tikzpicture}
\begin{axis}[width=10.5cm,height=5cm,xlabel={half-width (volatility points)},ylabel={profit an hour (per unit of vega)},
  tick label style={font=\small},label style={font=\small},
  xmin=0,xmax=1.5,ymin=-1,ymax=2.5,grid=major,grid style={gray!25},
  legend columns=5,legend cell align=left,legend style={font=\footnotesize,at={(0.5,-0.3)},anchor=north,draw=none,fill=none,/tikz/every even column/.append style={column sep=6pt}},
  table/col sep=comma]
\addplot[omDef,very thick] table[x=w,y=p0]{figdata/derivatives/26-options-market-making-in-practice/widths.csv};
\addplot[omThm,thick] table[x=w,y=p2]{figdata/derivatives/26-options-market-making-in-practice/widths.csv};
\addplot[omProp,thick] table[x=w,y=p5]{figdata/derivatives/26-options-market-making-in-practice/widths.csv};
\addplot[violet,thick] table[x=w,y=p10]{figdata/derivatives/26-options-market-making-in-practice/widths.csv};
\addplot[black,thick,dashed] table[x=w,y=p20]{figdata/derivatives/26-options-market-making-in-practice/widths.csv};
\addplot[black,only marks,mark=*,mark size=1.6pt] table[x=w,y=profit]{figdata/derivatives/26-options-market-making-in-practice/widths_opt.csv};
\legend{no informed,2 an hour,5,10,20}
\end{axis}
\end{tikzpicture}

\omcaption{Expected profit per hour of quoting a half-width around the theo, per unit of \omterm{def:dv:greeks-and-the-hedging-pnl:greeks}{vega}, for several rates of informed orders; dots mark the
best half-width. Informed flow pushes the width out and the profit down. Data: the tutorial.}\label{fig:dv:options-market-making-in-practice:width}
\end{omfigure}

Measured in volatility points, the best width does not depend on the option. In price it is that width times the option's \omterm{def:dv:greeks-and-the-hedging-pnl:greeks}{vega}. A one-month
at-the-money option on a share at 100 has a \omterm{def:dv:greeks-and-the-hedging-pnl:greeks}{vega} of 0.115 a point and a three-month option 0.199, so at 0.42 point the half-widths are 4.8 and 8.4 cents.
The inputs are measured, not guessed: the theo's error from mark-outs (Book 2, chapter 15) of past trades against later theos, and the informed
share from mark-outs by counterparty.

\section{Delta-hedging policies}

A market maker's book accumulates gamma from customers' trades, and its delta moves with every tick. Hedging each move costs the bid--ask spread of
the underlying; hedging seldom leaves risk. Time-based hedging, every so often, ignores where the delta is. A band hedges only when it must.

\begin{definition}[Hedging band]\label{def:dv:options-market-making-in-practice:band}
A \emph{hedging band}\index{hedging band} is a range around the target delta within which the hedge is left alone; when the delta leaves it, the
hedge trades back to the band's nearest edge, not to its centre, so that trades are only as large as the risk requires.
\end{definition}

The band's width follows from a balance. The delta diffuses at a rate set by the gamma, $\Gamma S\sigma$ per unit of $\sqrt t$, so it leaves a band of
half-width $H$ after a time of order $H^2/(\Gamma S\sigma)^2$ and costs about $\lambda SH$ in spread each time, with $\lambda$ the proportional cost.
The variance of the unhedged part grows as $H^2S^2\sigma^2$. Minimising the cost rate plus a risk-aversion $\gamma$ times the variance rate gives
$H^3\propto\lambda S\Gamma^2/\gamma$: a band half-width that scales as the cube root of cost times gamma squared. Whalley and Wilmott derived the optimal
band for small costs from a utility-maximising model as an expression in the option's gamma. The build uses $H=c\,(\lambda S\Gamma^2)^{1/3}$ and treats
$c$ as the desk's risk dial.

\begin{example}[Bands against the clock]\label{ex:dv:options-market-making-in-practice:bands}
A short one-month at-the-money straddle on a share at 100 (20\% volatility) is hedged over its life, checking 13 times a day, with a cost of 5 basis points
of the traded notional. Hedging at every check costs 0.525 on average and leaves a P\&L standard deviation of 0.321. Once a day costs 0.142 with a standard
deviation of 0.855. A band with $c=1$ costs 0.140, the same as daily hedging, with a standard deviation of 0.505, 41\% less. With $c=0.25$ it costs 0.313
with 0.313, and every-third-check hedging costs about the same (0.302) with 0.434
(\cref{fig:dv:options-market-making-in-practice:bands}).
\end{example}

\begin{omfigure}
\begin{tikzpicture}
\begin{axis}[width=10.5cm,height=5.2cm,xlabel={average hedging cost},ylabel={P\&L standard deviation},
  tick label style={font=\small},label style={font=\small},
  xmin=0,xmax=0.6,ymin=0,ymax=1.4,grid=major,grid style={gray!25},
  legend columns=2,legend cell align=left,legend style={font=\footnotesize,at={(0.5,-0.3)},anchor=north,draw=none,fill=none,/tikz/every even column/.append style={column sep=8pt}},
  table/col sep=comma]
\addplot[omProp,very thick,mark=*,mark size=1.6pt] table[x=cost,y=sd]{figdata/derivatives/26-options-market-making-in-practice/hedge_time.csv};
\addplot[omThm,very thick,mark=square*,mark size=1.6pt] table[x=cost,y=sd]{figdata/derivatives/26-options-market-making-in-practice/hedge_band.csv};
\legend{hedging on a clock (every 1 to 26 checks),hedging band ($c$ from 0.25 to 4)}
\end{axis}
\end{tikzpicture}

\omcaption{Hedging cost against the remaining risk for a short one-month straddle: time-based hedging at intervals from one to 26 checks (13 checks a day),
and bands of increasing width. The band's frontier lies below the clock's. Data: the tutorial.}\label{fig:dv:options-market-making-in-practice:bands}
\end{omfigure}

\section{Pin risk, dividends and early exercise}

Some of a market maker's risks are not in any Greek. Book 1, chapter 23 described pin risk. A desk short 1\,000 calls struck at 100 with the share closing
at 100.02 on expiry day does not know how many will be exercised, anywhere from none to all, so it does not know whether it will open Monday short
100\,000 shares or flat. It closes such positions before the bell or holds the share position the expected assignment implies and accepts the rest.

Dividends create the opposite kind of opportunity. On the eve of an ex-date, holders of deep in-the-money American calls should exercise when the dividend
exceeds the time value they give up, the value of the put with the same strike and expiry (chapter 6). Many do not. Pool, Stoll and Whaley found that in
1996--2006 more than half of the long positions that should have been exercised were not, at a cost to holders of over \$491 million, and that market
makers captured most of it.

\begin{definition}[Dividend play]\label{def:dv:options-market-making-in-practice:play}
A \emph{dividend play}\index{dividend play} is a trade on the eve of an ex-dividend date in which market makers buy and sell large, offsetting quantities of a
deep in-the-money call among themselves and exercise all their long calls; because assignment falls pro rata on all short positions, their short calls absorb
most of the non-assignment left by holders who fail to exercise, and each unassigned short call earns the dividend less the time value.
\end{definition}

The arithmetic is short. With $N$ public long contracts of which a share $f$ fails to exercise, and $q$ contracts traded twice between two market makers,
$2q+(1-f)N$ calls are exercised against $N+2q$ short positions. The market makers' unassigned short calls number $2q\,fN/(N+2q)$, which tends to $fN$ as
$q$ grows. Fees grow linearly in $q$, so there is a best size, $q^*=\bigl(\sqrt{2fNg\,N/c}-N\bigr)/2$, with $g$ the gain per unassigned call and $c$ the fees
per $q$.

\begin{example}[One call series]\label{ex:dv:options-market-making-in-practice:play}
A dividend of 50 cents (\$50 a contract), a put worth 2 cents (\$2), so a gain of \$48 per unassigned contract; 10\,000 public contracts; fees (an assumption)
of \$0.10 a contract per trade and \$0.05 per exercise, so \$0.50 per $q$ for four trades and two exercises. If 30\% of holders fail to exercise, the best trade
is 32\,947 contracts. It leaves 2\,605 short calls unassigned, 87\% of the 3\,000 that were not exercised, for a gross gain of \$125\,026, fees of \$16\,474 and a net
of \$108\,553. With 10\% failing the best trade is 16\,909 contracts and nets \$28\,591; with 50\%, 43\,990 contracts and \$193\,510
(\cref{fig:dv:options-market-making-in-practice:play}).
\end{example}

\begin{omfigure}
\begin{tikzpicture}
\begin{axis}[width=10.5cm,height=5cm,xlabel={contracts traded each way (thousands)},ylabel={expected net profit (\$ thousands)},
  tick label style={font=\small},label style={font=\small},
  xmin=0,xmax=100,ymin=-20,ymax=210,grid=major,grid style={gray!25},
  legend columns=3,legend cell align=left,legend style={font=\footnotesize,at={(0.5,-0.3)},anchor=north,draw=none,fill=none,/tikz/every even column/.append style={column sep=8pt}},
  table/col sep=comma]
\addplot[omDef,thick] table[x=q,y=f10]{figdata/derivatives/26-options-market-making-in-practice/dividend.csv};
\addplot[omProp,very thick] table[x=q,y=f30]{figdata/derivatives/26-options-market-making-in-practice/dividend.csv};
\addplot[omThm,thick] table[x=q,y=f50]{figdata/derivatives/26-options-market-making-in-practice/dividend.csv};
\addplot[black,only marks,mark=*,mark size=1.6pt] table[x=q,y=net]{figdata/derivatives/26-options-market-making-in-practice/dividend_opt.csv};
\legend{10\% fail to exercise,30\%,50\%}
\end{axis}
\end{tikzpicture}

\omcaption{Expected net profit of a \omterm{def:dv:options-market-making-in-practice:play}{dividend play} on one call series (10\,000 public contracts, \$48 gain per unassigned call, \$0.50 of fees per contract traded
each way) against the size of the trade, for three shares of holders failing to exercise; dots mark the best size. Data: the
tutorial.}\label{fig:dv:options-market-making-in-practice:play}
\end{omfigure}

The expected profit is not the realised one: assignment is random, and a desk that is assigned more than its share keeps less. The trade exists because
fees and assignment rules make it cheap, and it disappears where they do not.

\begin{dated}{2026-09}{Strategy fee caps}\label{dat:dv:options-market-making-in-practice:fees}
Cboe Options' fees schedule of 15 September 2026 caps market-maker and other non-customer transaction fees at \$0.00 for merger, short-stock-interest, reversal,
conversion and jelly-roll strategies executed in open outcry on the same day in the same class. Dividend strategies are not on that list, so a \omterm{def:dv:options-market-making-in-practice:play}{dividend play}
there pays ordinary fees.
\end{dated}

\section{Bucketed limits}

A market maker's \omterm{def:dv:greeks-and-the-hedging-pnl:greeks}{vega} is not one number. Long one-month and short one-year \omterm{def:dv:greeks-and-the-hedging-pnl:greeks}{vega} of the same size is flat in total and fully exposed to a change in the term
structure. Limits are set by bucket.

\begin{definition}[Vega bucket]\label{def:dv:options-market-making-in-practice:bucket}
A \emph{vega bucket}\index{vega bucket} is a range of expiries (and, in finer schemes, of moneyness) over which a book's \omterm{def:dv:greeks-and-the-hedging-pnl:greeks}{vega} is summed and limited; a book's
\omterm{def:dv:greeks-and-the-hedging-pnl:greeks}{vega} profile is its \omterm{def:dv:greeks-and-the-hedging-pnl:greeks}{vega} by bucket, and limits apply to each bucket separately.
\end{definition}

When a bucket approaches its limit, the quoting engine does not stop at once. It shades both quotes in the direction that attracts offsetting trades: a desk
short \omterm{def:dv:greeks-and-the-hedging-pnl:greeks}{vega} raises its bid and ask so that customers sell it options. It withdraws the side that would breach the limit only when the limit is reached.

\begin{example}[A month of put demand]\label{ex:dv:options-market-making-in-practice:toy}
A toy market maker quotes 15 series (one, three and six months; strikes 90 to 110) for 20 days, 13 checks a day. Orders of 10 contracts arrive three times a
check; uninformed customers buy puts two times in three and calls half the time; one order in five is informed. The theo errs by 0.5 point, the half-width is
0.4 point, and the delta is hedged at every check at 5 basis points. Limits are \$3\,000, \$5\,000 and \$6\,000 of \omterm{def:dv:greeks-and-the-hedging-pnl:greeks}{vega} per point. With shading, the book's six-month
\omterm{def:dv:greeks-and-the-hedging-pnl:greeks}{vega} peaks at $-\$3\,369$ and the P\&L is \$11\,141: edge \$36\,280, selection $-\$8\,720$, hedging $-\$9\,815$, inventory $-\$6\,605$. Without it, the same flow drives the
six-month \omterm{def:dv:greeks-and-the-hedging-pnl:greeks}{vega} to $-\$5\,033$, the inventory loses \$15\,240, and the month ends at $-\$607$ (\cref{fig:dv:options-market-making-in-practice:vega}). One path proves
nothing about averages. It shows where the risk goes when nothing leans against the flow.
\end{example}

\begin{omfigure}
\begin{tikzpicture}
\begin{groupplot}[group style={group size=2 by 1,horizontal sep=1.3cm},
  width=6.6cm,height=5cm,xlabel={check (13 a day)},xmin=0,xmax=260,ymin=-6.5,ymax=2.5,
  tick label style={font=\small},label style={font=\small},grid=major,grid style={gray!25},table/col sep=comma]
\nextgroupplot[ylabel={vega (\$ thousands a point)},title={\small three-month bucket},
  legend columns=2,legend cell align=left,legend style={font=\footnotesize,at={(1.08,-0.3)},anchor=north,draw=none,fill=none,/tikz/every even column/.append style={column sep=8pt}}]
\addplot[omThm,very thick] table[x=step,y=m3_on]{figdata/derivatives/26-options-market-making-in-practice/vega.csv};
\addplot[omProp,thick] table[x=step,y=m3_off]{figdata/derivatives/26-options-market-making-in-practice/vega.csv};
\addplot[gray,dashed,domain=0:260] {-5};
\legend{quotes shaded by the limits,no limits}
\nextgroupplot[title={\small six-month bucket},yticklabels={}]
\addplot[omThm,very thick] table[x=step,y=m6_on]{figdata/derivatives/26-options-market-making-in-practice/vega.csv};
\addplot[omProp,thick] table[x=step,y=m6_off]{figdata/derivatives/26-options-market-making-in-practice/vega.csv};
\addplot[gray,dashed,domain=0:260] {-6};
\end{groupplot}
\end{tikzpicture}

\omcaption{\omterm{def:dv:greeks-and-the-hedging-pnl:greeks}{Vega} by bucket of a toy options market maker facing a month of put demand, with quotes shaded as each bucket approaches its limit (dashed) and
without. Data: the tutorial.}\label{fig:dv:options-market-making-in-practice:vega}
\end{omfigure}

Gamma gets the same treatment by expiry, and the largest single names get their own buckets. The limits are only as good as the \omterm{def:dv:greeks-and-the-hedging-pnl:greeks}{Greeks} behind them, which is why
the desk's \omterm{def:dv:greeks-and-the-hedging-pnl:greeks}{Greeks} come from the same surface as its quotes.

\section{Tutorial: a toy options market maker}\label{tut:dv:options-market-making-in-practice:tut}

\begin{tutorial}
\textbf{Goal.} Compute the width model's best half-widths, compare \omterm{def:dv:options-market-making-in-practice:band}{hedging bands} with hedging on a clock, size a \omterm{def:dv:options-market-making-in-practice:play}{dividend play}, and run a toy market maker with
bucketed limits. \textbf{End state:} the four figures, the staleness table and the numbers of the weekend problem.
\begin{tutsteps}
\item \textbf{The hedger}, band or clock, vectorised over paths:
\omcode{code/firm/optmm/firm_optmm.py}{77}{108}{Delta hedging on a clock or in a band.}\label{lst:dv:options-market-making-in-practice:hedge}
\item \textbf{The \omterm{def:dv:options-market-making-in-practice:play}{dividend play}} and its best size:
\omcode{code/firm/optmm/firm_optmm.py}{118}{137}{The dividend play's expected profit.}\label{lst:dv:options-market-making-in-practice:play}
\item \textbf{The quote}, shaded by the bucket's use of its limit:
\omcode{code/firm/optmm/firm_optmm.py}{160}{174}{Quotes around the theo with limit shading.}\label{lst:dv:options-market-making-in-practice:quote}
\item \textbf{Run} \texttt{dv\_optmm.stale\_table()}, \texttt{width\_study()}, \texttt{hedge\_study()}, \texttt{dividend\_study()}, \texttt{toy\_mm()} and
\texttt{fig\_optmm.py}.
\end{tutsteps}
\textbf{What to change next.} Make the half-width depend on the informed share by series; hedge the toy market maker with a band; draw the \omterm{def:dv:options-market-making-in-practice:play}{dividend play}'s
assignment at random and plot the distribution of its profit.
\end{tutorial}

\section{Build: the options quoting engine}\label{bld:dv:options-market-making-in-practice:optmm}

\begin{build}
\bfield{Purpose} The miniature firm's options quoting engine: theos from a live surface, widths from a model, a band hedger, the \omterm{def:dv:options-market-making-in-practice:play}{dividend play}'s arithmetic,
and bucketed limits that shade quotes.
\bfield{Interface} \texttt{theo(option, spot, t, surface)} $\to$ (value, \omterm{def:dv:greeks-and-the-hedging-pnl:greeks}{vega} per point); \texttt{width\_profit}, \texttt{optimal\_width(err, lam\_u, w\_scale,
lam\_i)}; \texttt{band\_half\_width}, \texttt{hedge\_paths(paths, book, vol, dt, cost, band, every)}; \texttt{should\_exercise}, \texttt{dividend\_play(public,
fail, q, gain, trade\_fee, exercise\_fee)}, \texttt{best\_play}; \texttt{Bucket(name, lo, hi, limit)}, \texttt{vega\_by\_bucket}, \texttt{quote(option, spot, t,
surface, half\_width\_vol, bucket\_vega, limit)}.
\bfield{Rules} Widths are set in volatility points and converted by \omterm{def:dv:greeks-and-the-hedging-pnl:greeks}{vega}; hedges trade to the band's edge; a limit shades before it withdraws; every trade's
mark-out is recorded for the width model.
\bfield{Acceptance tests} \texttt{code/firm/optmm/tests/}: without informed flow the best width is the demand's scale, and informed flow widens it; an infinite
\omterm{def:dv:options-market-making-in-practice:play}{dividend play} captures every failure to exercise and the best size beats its neighbours; quotes are symmetric at zero use, shaded when long \omterm{def:dv:greeks-and-the-hedging-pnl:greeks}{vega}, and one side is
withdrawn at the limit; bands cost less than hedging at every check and hold more risk than tight bands.
\bfield{Stretch} Random assignment in the \omterm{def:dv:options-market-making-in-practice:play}{dividend play}; widths by series from recorded mark-outs; gamma limits by expiry.
\end{build}

\begin{omsources}
\item M.~Pool, H.~R.~Stoll and R.~E.~Whaley, ``Failure to exercise call options: an anomaly and a trading game'', \emph{Journal of Financial Markets} 11(1)
(2008) 1--35.
\item A.~E.~Whalley and P.~Wilmott, ``An asymptotic analysis of an optimal hedging model for option pricing with transaction costs'', \emph{Mathematical Finance}
7(3) (1997) 307--324.
\item Cboe Exchange, Inc., fees schedule, 15 September 2026, footnote 13.
\end{omsources}

\section{Exercises}

\begin{exercise}[$\star$]\label{exo:dv:options-market-making-in-practice:1}
Why does the best half-width equal the demand's scale (0.3 point) when there is no informed flow?
\end{exercise}

\begin{exercise}[$\star$]\label{exo:dv:options-market-making-in-practice:2}
Convert a half-width of 0.42 volatility point into cents for one-month and three-month at-the-money options on a share at 100.
\end{exercise}

\begin{exercise}[$\star$]\label{exo:dv:options-market-making-in-practice:3}
Check the \omterm{def:dv:options-market-making-in-practice:play}{dividend play}'s count of unassigned calls for $N=10\,000$, $f=0.3$, $q=32\,947$.
\end{exercise}

\begin{exercise}[$\star\star$]\label{exo:dv:options-market-making-in-practice:4}
Derive the best trade size $q^*$ of the \omterm{def:dv:options-market-making-in-practice:play}{dividend play}.
\end{exercise}

\begin{exercise}[$\star\star$]\label{exo:dv:options-market-making-in-practice:5}
Why does a \omterm{def:dv:options-market-making-in-practice:band}{hedging band} trade to its edge rather than to its centre?
\end{exercise}

\begin{exercise}[$\star\star$]\label{exo:dv:options-market-making-in-practice:6}
Why does the stale-surface error grow as the square root of the refit interval and fall with the option's expiry?
\end{exercise}

\begin{exercise}[$\star\star\star$]\label{exo:dv:options-market-making-in-practice:7}
\emph{Coding.} Draw the \omterm{def:dv:options-market-making-in-practice:play}{dividend play}'s assignments at random (each short call assigned with the same probability) for $f=0.3$ and the best size, and give the
standard deviation of the net profit.
\end{exercise}

\begin{exercise}[$\star\star\star$]\label{exo:dv:options-market-making-in-practice:8}
\emph{Find the flaw.} ``Our toy market maker made \$11\,141 with limits and lost \$607 without: limits add about \$11\,700 a month.''
\end{exercise}

\section{Problem: The Dividend-Play Morning}

\begin{problem}[{Weekend problem --- what the non-exercisers are worth}]\label{pb:dv:options-market-making-in-practice:1}
A share pays a 50-cent dividend tomorrow. Its deep in-the-money call series has 10\,000 contracts of public open interest; the put with the same strike is worth 2
cents. Two market makers consider a \omterm{def:dv:options-market-making-in-practice:play}{dividend play}.

\medskip\noindent\textbf{Part I --- Exercise.}
\begin{enumerate}
\item Should a holder of the call exercise tonight, and what does failing to cost per contract?
\item Why can early exercise of a call be optimal only just before an ex-date (with zero rates)?
\item What did Pool, Stoll and Whaley find about how often holders fail?
\item Who holds the other side of the public's calls, and what happens to them when holders fail?
\item Why does assignment pro rata over short positions make the play work?
\end{enumerate}

\medskip\noindent\textbf{Part II --- The play.}
\begin{enumerate}[resume]
\item Write the number of the market makers' unassigned short calls as a function of $q$.
\item What is the most they can capture, and at what size?
\item Give the fees per contract traded each way under the chapter's assumption.
\item Compute the best size and the net profit when 30\% fail.
\item Give the net profit when 10\% and when 50\% fail.
\end{enumerate}

\medskip\noindent\textbf{Part III --- Risks.}
\begin{enumerate}[resume]
\item What happens if the market makers are assigned more than their pro-rata share?
\item What happens if the public exercises more than expected?
\item How does the absence of a fee cap for dividend strategies change the trade?
\item What other risk does the desk carry overnight on these positions?
\item How would you estimate the failure rate for a series in advance?
\end{enumerate}

\medskip\noindent\textbf{Part IV --- Judgement.}
\begin{enumerate}[resume]
\item Who loses from the play, and is it the market makers' doing?
\item Should the desk do the trade at 10\% expected failure?
\item How would a rule change on exercise or fees end the play?
\item State the \emph{named result}: the expected profit of a \omterm{def:dv:options-market-making-in-practice:play}{dividend play} on one call series when a given share of holders fails to exercise, net of fees.
\item In one sentence: where does the \omterm{def:dv:options-market-making-in-practice:play}{dividend play}'s profit come from?
\end{enumerate}
\end{problem}

\section{Interview questions}

\begin{interviewq}[$\star$ \iqroles{trader}]\label{iq:dv:options-market-making-in-practice:1}
How do you set the width of an option quote?
\end{interviewq}

\begin{interviewq}[$\star\star$ \iqroles{trader, researcher}]\label{iq:dv:options-market-making-in-practice:2}
How often should an options market maker delta-hedge?
\end{interviewq}

\begin{interviewq}[$\star\star$ \iqroles{trader}]\label{iq:dv:options-market-making-in-practice:3}
Explain a \omterm{def:dv:options-market-making-in-practice:play}{dividend play}. Who loses?
\end{interviewq}

\begin{interviewq}[$\star\star$ \iqroles{developer}]\label{iq:dv:options-market-making-in-practice:4}
Design the part of a quoting engine that updates theos as the underlying ticks. What must be fast, and what can be slow?
\end{interviewq}

\begin{interviewq}[$\star\star$ \iqroles{risk, trader}]\label{iq:dv:options-market-making-in-practice:5}
Why limit \omterm{def:dv:greeks-and-the-hedging-pnl:greeks}{vega} by bucket rather than in total? What does the desk do as a bucket nears its limit?
\end{interviewq}

\begin{interviewq}[$\star\star\star$ \iqroles{trader}]\label{iq:dv:options-market-making-in-practice:6}
It is expiry Friday and the share is sitting on your largest short strike. What do you do?
\end{interviewq}
